\chapter{\texorpdfstring{Spin-$\frac12$ Algebra and Lattice Conventions}{Spin-1/2 Algebra and Lattice Conventions}}

\section{Purpose of this chapter}

Before discussing Jordan--Wigner transformations, Bogoliubov rotations, transfer
matrices, or Bethe ansatz, we must fix a small number of elementary conventions.
For spin chains, many sign errors do not come from difficult physics; they come
from inconsistent choices of
\begin{equation}
    \sigma^\pm,
    \qquad
    |\uparrow\rangle, |\downarrow\rangle,
    \qquad
    \sigma^z \text{ versus } n,
    \qquad
    \text{open versus periodic boundary conditions}.
\end{equation}
This chapter fixes these conventions once and for all. The main convention used
throughout the free-fermion part of the notes is
\begin{equation}
    |0\rangle_j\equiv |\uparrow\rangle_j,
    \qquad
    |1\rangle_j\equiv |\downarrow\rangle_j,
    \qquad
    n_j=\frac{1-\sigma_j^z}{2}.
\end{equation}
Thus the fermionic occupied state corresponds to a down spin. With this choice,
the Jordan--Wigner convention in the next chapter is
\begin{equation}
    c_j=
    \left(\prod_{\ell=1}^{j-1}\sigma_\ell^z\right)\sigma_j^+,
    \qquad
    c_j^\dagger=
    \left(\prod_{\ell=1}^{j-1}\sigma_\ell^z\right)\sigma_j^- .
\end{equation}
The present chapter should therefore be read as the local algebraic dictionary on
which all later fermionization formulae are based.

\section{Local Hilbert space}

A single spin-$1/2$ degree of freedom has local Hilbert space
\begin{equation}
    \calH_j\cong \CC^2.
\end{equation}
We use the ordered basis
\begin{equation}
    |\uparrow\rangle
    =
    \begin{pmatrix}1\\0\end{pmatrix},
    \qquad
    |\downarrow\rangle
    =
    \begin{pmatrix}0\\1\end{pmatrix}.
\end{equation}
The Pauli matrices are
\begin{equation}
    \sigma^x=
    \begin{pmatrix}
        0&1\\
        1&0
    \end{pmatrix},
    \qquad
    \sigma^y=
    \begin{pmatrix}
        0&-\ii\\
        \ii&0
    \end{pmatrix},
    \qquad
    \sigma^z=
    \begin{pmatrix}
        1&0\\
        0&-1
    \end{pmatrix}.
\end{equation}
They act on the basis states as
\begin{equation}
    \sigma^z|\uparrow\rangle=|\uparrow\rangle,
    \qquad
    \sigma^z|\downarrow\rangle=-|\downarrow\rangle.
\end{equation}
For a chain of length $L$, the many-body Hilbert space is the ordered tensor
product
\begin{equation}
    \calH_L
    =
    \bigotimes_{j=1}^{L}\calH_j
    \cong
    (\CC^2)^{\otimes L}.
\end{equation}
A basis vector is written as
\begin{equation}
    |s_1,s_2,\ldots,s_L\rangle,
    \qquad
    s_j\in\{\uparrow,\downarrow\}.
\end{equation}
Equivalently, with the occupation convention introduced above,
\begin{equation}
    |s_1,s_2,\ldots,s_L\rangle
    \equiv
    |n_1,n_2,\ldots,n_L\rangle,
    \qquad
    n_j=\begin{cases}
        0, & s_j=\uparrow,\\
        1, & s_j=\downarrow.
    \end{cases}
\end{equation}

The operator acting as $\sigma^\alpha$ on site $j$ and as the identity elsewhere is
\begin{equation}
    \sigma_j^\alpha
    =
    \id_1\otimes\cdots\otimes\id_{j-1}
    \otimes\sigma^\alpha
    \otimes\id_{j+1}\otimes\cdots\otimes\id_L,
    \qquad
    \alpha=x,y,z.
\end{equation}
Operators supported on distinct sites commute:
\begin{equation}
    [\sigma_i^\alpha,\sigma_j^\beta]=0,
    \qquad
    i\neq j.
\end{equation}

\section{Pauli algebra and spin operators}

On one site, the Pauli matrices obey
\begin{equation}
\label{eq:pauli-product-algebra}
    \sigma^\alpha\sigma^\beta
    =
    \delta_{\alpha\beta}\id
    +
    \ii\sum_{\gamma}\epsilon_{\alpha\beta\gamma}\sigma^\gamma,
    \qquad
    \alpha,\beta\in\{x,y,z\}.
\end{equation}
Equivalently,
\begin{equation}
    [\sigma^\alpha,\sigma^\beta]
    =
    2\ii\sum_\gamma\epsilon_{\alpha\beta\gamma}\sigma^\gamma,
    \qquad
    \{\sigma^\alpha,\sigma^\beta\}
    =
    2\delta_{\alpha\beta}\id .
\end{equation}
The physical spin operators are
\begin{equation}
    S^\alpha_j=\frac12\sigma^\alpha_j.
\end{equation}
Therefore
\begin{equation}
    [S_i^\alpha,S_j^\beta]
    =
    \ii\delta_{ij}
    \sum_\gamma\epsilon_{\alpha\beta\gamma}S_j^\gamma.
\end{equation}
The total spin components are
\begin{equation}
    S_{\rm tot}^\alpha
    =
    \sum_{j=1}^{L}S_j^\alpha
    =
    \frac12\sum_{j=1}^{L}\sigma_j^\alpha.
\end{equation}

\begin{remark}[Pauli normalization versus spin normalization]
Many condensed-matter texts write Heisenberg-type Hamiltonians using $S^\alpha$,
whereas the free-fermion spin-chain literature often writes them using Pauli
matrices. Since $S^\alpha=\sigma^\alpha/2$, one has
\begin{equation}
    \bm{S}_j\cdot\bm{S}_{j+1}
    =
    \frac14
    \left(
        \sigma_j^x\sigma_{j+1}^x
        +
        \sigma_j^y\sigma_{j+1}^y
        +
        \sigma_j^z\sigma_{j+1}^z
    \right).
\end{equation}
Thus coupling constants differ by factors of $4$ when translating between the two
normalizations. Unless explicitly stated otherwise, the spin-chain Hamiltonians in
these notes are written in Pauli normalization.
\end{remark}

Useful trace identities are
\begin{equation}
    \Tr\,\sigma^\alpha=0,
    \qquad
    \Tr(\sigma^\alpha\sigma^\beta)=2\delta_{\alpha\beta},
\end{equation}
and
\begin{equation}
    \Tr(\sigma^\alpha\sigma^\beta\sigma^\gamma)
    =
    2\ii\epsilon_{\alpha\beta\gamma}.
\end{equation}
These identities are especially useful when deriving transfer matrices, density
matrices, and two-band Bloch Hamiltonians.

\section{Raising and lowering operators}

We define
\begin{equation}
\label{eq:sigma-pm-definition}
    \sigma^+
    =
    \frac{\sigma^x+\ii\sigma^y}{2}
    =
    \begin{pmatrix}
        0&1\\
        0&0
    \end{pmatrix},
    \qquad
    \sigma^-
    =
    \frac{\sigma^x-\ii\sigma^y}{2}
    =
    \begin{pmatrix}
        0&0\\
        1&0
    \end{pmatrix}.
\end{equation}
Hence
\begin{equation}
    \sigma^+|\downarrow\rangle=|\uparrow\rangle,
    \qquad
    \sigma^-|\uparrow\rangle=|\downarrow\rangle,
\end{equation}
while
\begin{equation}
    \sigma^+|\uparrow\rangle=0,
    \qquad
    \sigma^-|\downarrow\rangle=0.
\end{equation}
The inverse relations are
\begin{equation}
\label{eq:sigma-x-y-pm}
    \sigma^x=\sigma^++\sigma^-,
    \qquad
    \sigma^y=\ii(\sigma^- -\sigma^+).
\end{equation}
The local algebra is
\begin{equation}
    (\sigma^+)^2=(\sigma^-)^2=0,
    \qquad
    \{\sigma^+,\sigma^-\}=\id,
    \qquad
    [\sigma^+,\sigma^-]=\sigma^z,
\end{equation}
and
\begin{equation}
    [\sigma^z,\sigma^+]=2\sigma^+,
    \qquad
    [\sigma^z,\sigma^-]=-2\sigma^-.
\end{equation}
The two local projectors are
\begin{equation}
\label{eq:spin-projectors}
    \sigma^+\sigma^-=\frac{1+\sigma^z}{2}=|\uparrow\rangle\langle\uparrow|,
    \qquad
    \sigma^-\sigma^+=\frac{1-\sigma^z}{2}=|\downarrow\rangle\langle\downarrow|.
\end{equation}

\section{Occupation convention and hard-core bosons}

The Jordan--Wigner transformation used later identifies the spin-down state with
fermion occupation. We therefore define
\begin{equation}
\label{eq:spin-occupation-convention}
    |0\rangle_j\equiv |\uparrow\rangle_j,
    \qquad
    |1\rangle_j\equiv |\downarrow\rangle_j.
\end{equation}
The local number operator is
\begin{equation}
\label{eq:number-spin-convention}
    n_j
    =
    |\downarrow\rangle_j\langle\downarrow|
    =
    \sigma_j^-\sigma_j^+
    =
    \frac{1-\sigma_j^z}{2}.
\end{equation}
Thus
\begin{equation}
\label{eq:sigmaz-number-convention}
    \sigma_j^z=1-2n_j.
\end{equation}

It is convenient to introduce local hard-core boson operators
\begin{equation}
\label{eq:hard-core-boson-convention}
    b_j=\sigma_j^+,
    \qquad
    b_j^\dagger=\sigma_j^-.
\end{equation}
With the present convention, $b_j^\dagger$ creates a down spin, hence creates an
occupied site:
\begin{equation}
    b_j^\dagger |0\rangle_j=|1\rangle_j,
    \qquad
    b_j |1\rangle_j=|0\rangle_j.
\end{equation}
The number operator is
\begin{equation}
    n_j=b_j^\dagger b_j.
\end{equation}
On the same site, hard-core bosons obey
\begin{equation}
    b_j^2=(b_j^\dagger)^2=0,
    \qquad
    \{b_j,b_j^\dagger\}=1,
    \qquad
    [b_j,b_j^\dagger]=1-2n_j=\sigma_j^z.
\end{equation}
On different sites, however, they commute rather than anticommute:
\begin{equation}
    [b_i,b_j]=0,
    \qquad
    [b_i,b_j^\dagger]=0,
    \qquad
    [b_i^\dagger,b_j^\dagger]=0,
    \qquad
    i\neq j.
\end{equation}
This mixed behavior is the reason a nonlocal Jordan--Wigner string is needed to
convert hard-core bosons into canonical fermions.

\begin{remark}[Why $b=\sigma^+$ but $b^\dagger=\sigma^-$]
This convention may look reversed if one thinks of $\sigma^+$ as a creation
operator. Here the created object is not spin-up, but an occupied fermionic site.
Since occupation means spin down, the operator that creates occupation must send
$|\uparrow\rangle$ to $|\downarrow\rangle$. Therefore $b^\dagger=\sigma^-$.
\end{remark}

\section{Useful two-site identities}

The following identities are used repeatedly in the XY, XX, XXZ, and Kitaev-chain
chapters. For nearest-neighbour sites $j$ and $j+1$,
\begin{align}
\label{eq:xy-raising-lowering-identity}
    \sigma_j^x\sigma_{j+1}^x+
    \sigma_j^y\sigma_{j+1}^y
    &=
    2\left(
        \sigma_j^+\sigma_{j+1}^-
        +
        \sigma_j^-\sigma_{j+1}^+
    \right),
    \\
\label{eq:pairing-raising-lowering-identity}
    \sigma_j^x\sigma_{j+1}^x-
    \sigma_j^y\sigma_{j+1}^y
    &=
    2\left(
        \sigma_j^+\sigma_{j+1}^+
        +
        \sigma_j^-\sigma_{j+1}^-
    \right).
\end{align}
Equivalently,
\begin{align}
    \sigma_j^x\sigma_{j+1}^x
    &=
    \sigma_j^+\sigma_{j+1}^+
    +
    \sigma_j^+\sigma_{j+1}^-
    +
    \sigma_j^-\sigma_{j+1}^+
    +
    \sigma_j^-\sigma_{j+1}^-,
    \\
    \sigma_j^y\sigma_{j+1}^y
    &=
    -\sigma_j^+\sigma_{j+1}^+
    +
    \sigma_j^+\sigma_{j+1}^-
    +
    \sigma_j^-\sigma_{j+1}^+
    -
    \sigma_j^-\sigma_{j+1}^- .
\end{align}
In terms of the occupation convention,
\begin{equation}
\label{eq:zz-number-identity}
    \sigma_j^z\sigma_{j+1}^z
    =
    (1-2n_j)(1-2n_{j+1})
    =
    1-2n_j-2n_{j+1}+4n_jn_{j+1}.
\end{equation}
This last identity explains why an Ising interaction in the $z$ direction becomes a
density--density interaction after fermionization, whereas the $xx$ and $yy$ terms
of the XY chain become hopping and pairing terms.

\section{Lattice conventions}

Unless stated otherwise, a one-dimensional chain has sites
\begin{equation}
    j=1,2,\ldots,L,
\end{equation}
with lattice spacing set to
\begin{equation}
    a=1.
\end{equation}
A nearest-neighbour Hamiltonian is written using one of two bond sets:
\begin{align}
    \mathcal{B}_{\rm OBC}
    &=
    \{(j,j+1)\,|\,j=1,2,\ldots,L-1\},
    \\
    \mathcal{B}_{\rm PBC}
    &=
    \{(j,j+1)\,|\,j=1,2,\ldots,L\},
    \qquad
    L+1\equiv 1.
\end{align}
Thus
\begin{equation}
    \sum_{\langle j,j+1\rangle\in\mathcal{B}_{\rm OBC}}
    \quad\text{means}\quad
    \sum_{j=1}^{L-1},
\end{equation}
and
\begin{equation}
    \sum_{\langle j,j+1\rangle\in\mathcal{B}_{\rm PBC}}
    \quad\text{means}\quad
    \sum_{j=1}^{L}\quad\text{with}\quad L+1\equiv1.
\end{equation}
When no boundary condition is explicitly stated, formulae for local bulk
Hamiltonian densities should be read as open-chain formulae first. Periodic
boundary conditions must be treated with additional care after the Jordan--Wigner
transformation because the boundary spin term carries a fermion-parity-dependent
sign.

For translation-invariant periodic chains, the one-site translation operator $T$
is defined by
\begin{equation}
    T\sigma_j^\alpha T^{-1}=\sigma_{j+1}^\alpha,
    \qquad
    T^L=1.
\end{equation}
Momentum eigenstates satisfy
\begin{equation}
    T|\psi_k\rangle=\ee^{\ii k}|\psi_k\rangle,
\end{equation}
where the allowed values of $k$ depend on the boundary condition and, after the
Jordan--Wigner transformation, on the fermion-parity sector. The detailed Fourier
convention is fixed in the Fourier-transform chapter.

\begin{remark}[Ordering convention]
The Jordan--Wigner string depends on a one-dimensional ordering of sites. In these
notes the ordering is always
\begin{equation}
    1<2<\cdots<L.
\end{equation}
This is natural for one-dimensional chains. In higher-dimensional spin models, a
linear ordering can still be imposed, but it usually destroys locality. This is one
reason why two-dimensional exact solvability often requires other structures, such
as commuting projectors, transfer matrices, or Majorana fermions coupled to static
$\ZZ_2$ gauge fields.
\end{remark}

\section{Common spin Hamiltonians}

We now list the spin-chain Hamiltonians that will appear repeatedly in these
notes. The purpose here is not to solve them, but to fix signs and normalizations.
The default Pauli-normalized nearest-neighbour XYZ Hamiltonian is
\begin{equation}
\label{eq:xyz-general-convention}
    H_{\rm XYZ}
    =
    -\sum_{\langle j,j+1\rangle}
    \left(
        J_x\sigma_j^x\sigma_{j+1}^x
        +
        J_y\sigma_j^y\sigma_{j+1}^y
        +
        J_z\sigma_j^z\sigma_{j+1}^z
    \right)
    -
    h\sum_{j=1}^{L}\sigma_j^z .
\end{equation}
With this convention, positive $J_\alpha$ favors ferromagnetic alignment in the
$\alpha$ direction.

\subsection{Ising and transverse-field Ising chains}

The Ising interaction used in the free-fermion part of the notes is chosen along
$x$, while the transverse field is along $z$:
\begin{equation}
\label{eq:tfim-convention-chapter}
    H_{\rm TFIM}
    =
    -J\sum_{\langle j,j+1\rangle}\sigma_j^x\sigma_{j+1}^x
    -
    h\sum_{j=1}^{L}\sigma_j^z .
\end{equation}
The choice of $x$ as the ordering direction and $z$ as the transverse-field
direction is adapted to the occupation convention $n_j=(1-\sigma_j^z)/2$. It is
unitarily equivalent to the alternative textbook convention
$-J\sum_j\sigma_j^z\sigma_{j+1}^z-h\sum_j\sigma_j^x$ by a global spin rotation.

The purely classical Ising interaction is recovered by setting $h=0$:
\begin{equation}
    H_{\rm Ising}^{x}
    =
    -J\sum_{\langle j,j+1\rangle}\sigma_j^x\sigma_{j+1}^x .
\end{equation}
A longitudinal field along the ordering direction,
\begin{equation}
    -h_x\sum_j\sigma_j^x,
\end{equation}
can be added, but it destroys the standard free-fermion solvability of the
one-dimensional TFIM.

\subsection{XY and XX chains}

The transverse-field XY chain is the central free-fermion mother model of these
notes:
\begin{equation}
\label{eq:xy-convention-chapter}
    H_{\rm XY}
    =
    -\sum_{\langle j,j+1\rangle}
    \left[
        \frac{J(1+\gamma)}{2}\sigma_j^x\sigma_{j+1}^x
        +
        \frac{J(1-\gamma)}{2}\sigma_j^y\sigma_{j+1}^y
    \right]
    -
    h\sum_{j=1}^{L}\sigma_j^z .
\end{equation}
The parameter $\gamma$ is the XY anisotropy. Important limits are
\begin{align}
    \gamma=1
    &\quad\Longrightarrow\quad
    H_{\rm XY}=H_{\rm TFIM},
    \\
    \gamma=0
    &\quad\Longrightarrow\quad
    H_{\rm XX}
    =
    -\frac{J}{2}
    \sum_{\langle j,j+1\rangle}
    \left(
        \sigma_j^x\sigma_{j+1}^x
        +
        \sigma_j^y\sigma_{j+1}^y
    \right)
    -h\sum_j\sigma_j^z .
\end{align}
Using \cref{eq:xy-raising-lowering-identity}, the XX interaction can also be
written as
\begin{equation}
    H_{\rm XX}
    =
    -J\sum_{\langle j,j+1\rangle}
    \left(
        \sigma_j^+\sigma_{j+1}^-
        +
        \sigma_j^-\sigma_{j+1}^+
    \right)
    -h\sum_j\sigma_j^z .
\end{equation}
This form makes the conservation of total $z$-magnetization manifest.

\subsection{XXZ and Heisenberg chains}

In the Pauli normalization compatible with the above XX convention, we write
\begin{equation}
\label{eq:xxz-pauli-convention}
    H_{\rm XXZ}
    =
    -\frac{J}{2}
    \sum_{\langle j,j+1\rangle}
    \left(
        \sigma_j^x\sigma_{j+1}^x
        +
        \sigma_j^y\sigma_{j+1}^y
        +
        \Delta\sigma_j^z\sigma_{j+1}^z
    \right)
    -h\sum_j\sigma_j^z .
\end{equation}
The isotropic point $\Delta=1$ gives a Heisenberg chain in Pauli normalization:
\begin{equation}
    H_{\rm Heis}
    =
    -\frac{J}{2}
    \sum_{\langle j,j+1\rangle}
    \bm{\sigma}_j\cdot\bm{\sigma}_{j+1}
    -h\sum_j\sigma_j^z .
\end{equation}
In spin-operator normalization, the same interaction is written as
\begin{equation}
    H_{\rm Heis}^{S}
    =
    \mathcal{J}\sum_{\langle j,j+1\rangle}
    \bm{S}_j\cdot\bm{S}_{j+1}
    -h_S\sum_j S_j^z.
\end{equation}
The conversion is
\begin{equation}
    \bm{S}_j\cdot\bm{S}_{j+1}
    =
    \frac14\bm{\sigma}_j\cdot\bm{\sigma}_{j+1},
    \qquad
    S_j^z=\frac12\sigma_j^z.
\end{equation}
The Bethe-ansatz chapters may use the conventional antiferromagnetic sign
$\mathcal{J}>0$; when this happens, the sign and normalization will be stated
explicitly.

\subsection{Bond-dependent Kitaev-type spin Hamiltonians}

In two dimensions, the Kitaev honeycomb model uses bond-dependent Ising
interactions:
\begin{equation}
\label{eq:kitaev-honeycomb-convention}
    H_{\rm Kitaev}
    =
    -J_x\sum_{\langle ij\rangle_x}\sigma_i^x\sigma_j^x
    -J_y\sum_{\langle ij\rangle_y}\sigma_i^y\sigma_j^y
    -J_z\sum_{\langle ij\rangle_z}\sigma_i^z\sigma_j^z.
\end{equation}
Here $\langle ij\rangle_\alpha$ denotes a nearest-neighbour bond of type $\alpha$.
This model is not solved by the one-dimensional Jordan--Wigner transformation.
Instead, it becomes exactly solvable through a Majorana representation and static
$\ZZ_2$ gauge fields.

\section{Symmetries and conserved quantities}

Symmetry is one of the main tools for block diagonalization. We collect the most
important symmetries here.

\subsection{Translation symmetry}

For periodic boundary conditions, a translation-invariant Hamiltonian satisfies
\begin{equation}
    [H,T]=0.
\end{equation}
The Hilbert space can then be decomposed into momentum sectors labelled by the
eigenvalues $\ee^{\ii k}$ of $T$. Translation symmetry is essential for the Fourier
transform and the reduction of quadratic Hamiltonians into independent momentum
blocks.

\subsection{Fermion parity and Ising spin flip}

Define the global $z$-parity operator
\begin{equation}
\label{eq:global-z-parity}
    P_z
    =
    \prod_{j=1}^{L}\sigma_j^z.
\end{equation}
It satisfies
\begin{equation}
    P_z^2=1,
    \qquad
    P_z\sigma_j^xP_z=-\sigma_j^x,
    \qquad
    P_z\sigma_j^yP_z=-\sigma_j^y,
    \qquad
    P_z\sigma_j^zP_z=\sigma_j^z.
\end{equation}
Therefore the TFIM and the XY chain commute with $P_z$:
\begin{equation}
    [H_{\rm TFIM},P_z]=0,
    \qquad
    [H_{\rm XY},P_z]=0.
\end{equation}
Under the occupation convention \cref{eq:number-spin-convention},
\begin{equation}
    P_z
    =
    \prod_{j=1}^{L}(1-2n_j)
    =
    (-1)^N,
    \qquad
    N=\sum_{j=1}^{L}n_j.
\end{equation}
Thus the global Ising spin-flip symmetry of the spin chain becomes fermion parity
after the Jordan--Wigner transformation.

\subsection{\texorpdfstring{$U(1)$}{U(1)} spin-rotation symmetry}

The total $z$-magnetization is
\begin{equation}
    M_z=\sum_{j=1}^{L}\sigma_j^z=L-2N.
\end{equation}
A Hamiltonian has $U(1)$ spin-rotation symmetry about the $z$ axis if
\begin{equation}
    [H,M_z]=0.
\end{equation}
The XX, XXZ, and Heisenberg chains have this symmetry. Equivalently, they are
invariant under
\begin{equation}
    U_z(\theta)
    =
    \exp\!\left(-\frac{\ii\theta}{2}\sum_j\sigma_j^z\right),
\end{equation}
which acts as
\begin{equation}
    U_z(\theta)\sigma_j^\pm U_z(\theta)^{-1}
    =
    \ee^{\mp\ii\theta}\sigma_j^\pm.
\end{equation}
The anisotropic XY chain with $\gamma\neq0$ does not conserve $M_z$ because
\cref{eq:pairing-raising-lowering-identity} contains pair creation and pair
annihilation terms. Nevertheless, it still conserves $P_z=(-1)^N$.

\subsection{Full spin-rotation symmetry}

The isotropic Heisenberg chain at zero field has full $SU(2)$ symmetry:
\begin{equation}
    [H_{\rm Heis},S_{\rm tot}^\alpha]=0,
    \qquad
    \alpha=x,y,z.
\end{equation}
An external field along $z$ breaks this to the $U(1)$ subgroup generated by
$S_{\rm tot}^z$.

\subsection{Boundary-sector subtlety after Jordan--Wigner}

For an open chain, the Jordan--Wigner transformation maps local nearest-neighbour
bulk terms to local fermion bilinears without any boundary ambiguity. For a
periodic spin chain, the boundary term connecting $L$ and $1$ depends on the parity
sector
\begin{equation}
    P_z=(-1)^N=p,
    \qquad
    p=\pm1.
\end{equation}
With the Jordan--Wigner convention used in these notes, the effective fermionic
boundary condition in a fixed parity sector is
\begin{equation}
\label{eq:chapter-effective-boundary-condition}
    c_{L+1}=-p\,c_1.
\end{equation}
Thus
\begin{equation}
    p=+1
    \quad\Longrightarrow\quad
    c_{L+1}=-c_1,
    \qquad
    p=-1
    \quad\Longrightarrow\quad
    c_{L+1}=c_1.
\end{equation}
That is, even fermion parity corresponds to anti-periodic fermions, while odd
fermion parity corresponds to periodic fermions. The derivation is given in the
Jordan--Wigner chapter.

\section{Summary of conventions}

The conventions fixed in this chapter are summarized as follows:
\begin{equation}
    |0\rangle_j=|\uparrow\rangle_j,
    \qquad
    |1\rangle_j=|\downarrow\rangle_j,
    \qquad
    n_j=\frac{1-\sigma_j^z}{2}.
\end{equation}
The raising and lowering operators are
\begin{equation}
    \sigma^+=\frac{\sigma^x+\ii\sigma^y}{2},
    \qquad
    \sigma^-=\frac{\sigma^x-\ii\sigma^y}{2},
\end{equation}
with
\begin{equation}
    \sigma^+|\downarrow\rangle=|\uparrow\rangle,
    \qquad
    \sigma^-|\uparrow\rangle=|\downarrow\rangle.
\end{equation}
The hard-core boson convention is
\begin{equation}
    b_j=\sigma_j^+,
    \qquad
    b_j^\dagger=\sigma_j^-,
    \qquad
    n_j=b_j^\dagger b_j.
\end{equation}
The principal free-fermion mother model is the transverse-field XY chain
\begin{equation}
    H_{\rm XY}
    =
    -\sum_{\langle j,j+1\rangle}
    \left[
        \frac{J(1+\gamma)}{2}\sigma_j^x\sigma_{j+1}^x
        +
        \frac{J(1-\gamma)}{2}\sigma_j^y\sigma_{j+1}^y
    \right]
    -h\sum_j\sigma_j^z,
\end{equation}
which contains the TFIM at $\gamma=1$ and the XX chain at $\gamma=0$. The global
operator
\begin{equation}
    P_z=\prod_j\sigma_j^z=(-1)^N
\end{equation}
will become fermion parity in the Jordan--Wigner representation. These conventions
are the reference point for all later equivalence maps.
